\section{总结与延伸}

\subsection{核心结论}

\begin{enumerate}
\item AI for Math 不等于刷数学题，而是从问题、直觉、形式化到验证的完整系统。
\item Lean/Coq 等 proof assistant 为 AI 提供硬反馈，使数学成为适合训练和评估的高价值领域。
\item 数学直觉和形式验证不是对立关系，强系统需要同时拥有 free attention 和 bounded attention。
\item Axiom 这类 Neo Lab 需要数学家网络、形式化系统、模型训练、资金算力和产品目标共同成立。
\item AI for Math 的终点不是替代数学家，而是与数学家共建证明、定理库和研究路线。
\end{enumerate}

\subsection{开放问题}

\begin{itemize}
\item AI 能否生成真正有审美和解释力的证明，而不只是长而正确的证明？
\item Lean 形式化会成为数学研究主流接口，还是只服务部分领域？
\item 数学家反馈如何规模化，才能形成 AI for Math 数据飞轮？
\item Axiom 这类 Neo Lab 的产品形态会是 proof assistant、研究平台，还是数学家协作系统？
\end{itemize}

\subsection{拓展阅读}

\begin{itemize}
\item Lean theorem prover 与 mathlib：理解形式化数学的基础设施。
\item Terence Tao、Kevin Buzzard 等关于 formalization 的公开讨论：理解数学社区为什么重视 Lean。
\item AI theorem proving、autoformalization、proof search 相关论文：理解 AI for Math 的技术路线。
\item EP139 和 EP138 笔记：把 AI for Math 放回 Agent 与模型实验室的更大图谱中。
\end{itemize}

\begin{importantbox}{最后的判断}
AI for Math 的难点在于，它同时要求模型懂数学、会写形式语言、能接受机器验证、也能与数学家直觉协作。它是 AI 从“生产力工具”走向“科学发现伙伴”的关键试验场。
\end{importantbox}

\end{document}
